\appendix
\section{Query Formulations}
\label{sec:appendix-queries}

% Revised 28 Oct 2026 against scripts/pathfinding_benchmark.py and scripts/tools/update_concurrent_driver.cpp at the
% frozen commit 7e9778c0. The text is what the harness runs; the paper's description of the baselines must agree with it.

Every number in Section~\ref{sec:eval} comes from the query text below, which is the fastest formulation we found for
each system; the harness and the result archive are released with the paper. Conventions: \texttt{\$name} is a
parameter bound by the driver, \texttt{\{name\}} a literal interpolated into the text. Every query returns the count of
reached pairs and the sum of the hop lengths. For the native systems self pairs are removed before the query and
counted client-side (and in sources $\times$ all the sources themselves are added); \duckpgq receives the full pair
table and answers a self pair with length 0. The comparisons below that justify a formulation were measured while the
harness was developed, not in the frozen run. Every Cypher pattern carries an upper bound of 30 hops, Kuzu's default
maximum; the GDS procedures and the \duckpgq queries carry none.

\paragraph{\duckpgq.} The sweeps drive the operator through the marker function that the \texttt{MATCH} rewrite of
Section~\ref{sec:planner} produces, so that the pairs come from a table and not from the cross product of two vertex
patterns. The equivalent \sqlpgq for one pair is
\begin{lstlisting}[language=SQL]
SELECT path_length(p) FROM GRAPH_TABLE (snb MATCH p = ANY SHORTEST
  (a:Person WHERE a.id = {src})-[k:Knows]->*(b:Person WHERE b.id = {dst}) COLUMNS (p));
\end{lstlisting}
Every run is a fresh \duckdb process with \texttt{SET experimental\_path\_finding\_operator=true}; the database is
attached read-only. The build route runs
\begin{lstlisting}[language=SQL]
SELECT count(len), coalesce(sum(len), 0) FROM (SELECT iterativelengthoperator(src, dst,
    struct_pack(src := pathfinding_edge_src, dst := pathfinding_edge_dst),
    {V}::BIGINT, {E}::BIGINT, 'graph:sweep') AS len
  FROM ldbc.benchmark_pairs_{n} AS pathfinding_pairs,
    (SELECT a.rowid::BIGINT AS pathfinding_edge_src, c.rowid::BIGINT AS pathfinding_edge_dst
     FROM ldbc.person_knows_person k JOIN ldbc.person a ON a.id = k.person1id
                                     JOIN ldbc.person c ON c.id = k.person2id) pathfinding_edges);
\end{lstlisting}
The index route creates \texttt{CREATE INDEX IF NOT EXISTS pathfinding\_benchmark\_csr ON ldbc.person\_knows\_person
USING CSR (person1id, person2id)} once, untimed, materializes the key pairs untimed into
\texttt{pathfinding\_index\_pairs} (the pair table joined to \texttt{person} on row id), and runs
\begin{lstlisting}[language=SQL]
SELECT count(len), coalesce(sum(len), 0) FROM (SELECT iterativelengthoperator(src, dst,
    0::BIGINT, 0::BIGINT, 'csr-index|ldbc|main|person_knows_person|pathfinding_benchmark_csr') AS len
  FROM pathfinding_index_pairs);
\end{lstlisting}
Sources $\times$ all uses \texttt{reachoperator} in place of \texttt{iterativelengthoperator}, with each source as both
source and destination, and sums the length of every reached vertex. The LDBC queries are undirected. On the build
route they use \texttt{person\_knows\_person\_sym}, the union of both directions, and the key
\texttt{'graph:sweep|undirected'}, which tells the operator that the edge set is symmetric; on the index route they use
an index created \texttt{WITH (undirected = true)}. In each process the timed statement is preceded by one untimed run
of the same statement, which also loads the index from disk.

\paragraph{Kuzu 0.11.3.} Embedded, on a read-only database, with 32 threads. One query per pair, with literal ids; a
pair table bound at both ends ran two orders of magnitude slower (15\,s against 0.15\,s for 64 pairs at SF1), and a
source bound through \texttt{WITH} made the recursive join start from every vertex.
\begin{lstlisting}
MATCH (a:Person {id: {src}})-[e:Knows* SHORTEST 1..30]->(b:Person {id: {dst}}) RETURN length(e)
MATCH (a:Person {id: {source}})-[e:Knows* SHORTEST 1..30]->(b:Person) RETURN count(*), sum(length(e))
\end{lstlisting}
The 32-connection variant opens 32 connections on one database with $\mathit{threads}/32$ threads each and sends each
connection its share of the pairs, assigned round-robin.

\paragraph{Neo4j 5.26.32 Enterprise.} One query per pair list with the parallel runtime:
\begin{lstlisting}
CYPHER runtime=parallel UNWIND $pairs AS pr MATCH (a:Person {id: pr[0]}), (b:Person {id: pr[1]})
  MATCH p = shortestPath((a)-[:KNOWS*..30]->(b)) RETURN count(p), sum(length(p))
CYPHER runtime=parallel UNWIND $srcs AS sid MATCH (a:Person {id: sid}) MATCH (b:Person) WHERE a <> b
  MATCH p = shortestPath((a)-[:KNOWS*..30]->(b)) RETURN count(p), sum(length(p))
\end{lstlisting}
The GDS rows use a projection built once after loading and the library's unlicensed concurrency of four:
\texttt{gds.shortestPath.dijkstra.stream} per pair and \texttt{gds.allShortestPaths.delta.stream} per source. GDS
cannot run under the parallel runtime and is absent from the update workload because its projection is static.

\paragraph{Memgraph 3.14.0.} Bounded BFS, one thread per query. The \texttt{WITH a} barrier is required; without it
Memgraph scans every $b$ and filters $a$ after the search (35\,s against 10\,ms at SF1). BI19 passes its source ids as
a list for the same reason.
\begin{lstlisting}
UNWIND $pairs AS pr MATCH (a:Person {id: pr[0]}), (b:Person {id: pr[1]})
  MATCH p = (a)-[:KNOWS *BFS ..30]->(b) RETURN count(p), sum(size(p))
MATCH (a:Person) WHERE a.id IN $srcs WITH a MATCH p = (a)-[:KNOWS *BFS ..30]->(b:Person)
  RETURN count(p), sum(size(p))
\end{lstlisting}

\paragraph{FalkorDB 4.2.2.} One thread per query. \texttt{shortestPath} is allowed only in \texttt{WITH} or
\texttt{RETURN}; there is no single-source form, so sources $\times$ all is one \texttt{shortestPath} per destination,
and no undirected form, so the LDBC queries run on a graph with every edge inserted in both directions.
\begin{lstlisting}
UNWIND $pairs AS pr MATCH (a:Person {id: pr[0]}), (b:Person {id: pr[1]})
  WITH shortestPath((a)-[:KNOWS*..30]->(b)) AS p RETURN count(p), sum(length(p))
MATCH (a:Person {id: $src}), (b:Person) WITH shortestPath((a)-[:KNOWS*..30]->(b)) AS p RETURN count(p), sum(length(p))
\end{lstlisting}

\paragraph{LDBC queries.} We run hop-count variants of IC13, BI19 and BI20 with five parameter sets each, chosen by the
harness: IC13 takes its pairs from the random pair table, BI19 pairs the cities ranked $i$ and $i+5$ by population,
and BI20 takes the five largest companies, each with a random person. Neo4j, Memgraph and Kuzu use undirected patterns
(\texttt{-[:KNOWS*..30]-}, \texttt{-[e:Knows* SHORTEST 1..30]-}), GDS an undirected projection, and FalkorDB its
directed pattern over the doubled edges. Kuzu runs BI19 as one single-source query per person of the first city with
the destinations filtered by city, and BI20 as one query per (employee, person) pair. With 32 client connections the
source ids are fetched first, inside the timed region, and split over the connections.

\paragraph{Update workload.} The LDBC insert and delete stream for \texttt{Person\_knows\_Person} is restricted to the
persons of the snapshot, with duplicate inserts and unmatched or repeated deletes dropped, and applied in batches of
2000 operations. \duckdb applies a batch as one transaction, \texttt{INSERT ... SELECT} from the staged log followed by
\texttt{DELETE ... USING}. Neo4j and Memgraph run the inserts and the deletes as two \texttt{UNWIND \$ops} statements
in one explicit transaction; Kuzu reads the batch from an \texttt{Op} node table and uses two transactions, because a
\texttt{MATCH ... DELETE} in the transaction that created the relationship did not see it; FalkorDB runs two
autocommit queries. After each batch every system answers the single-source query to the same 1024 destinations.
Memgraph and FalkorDB filter with \texttt{WHERE b.id IN \$dsts}; Neo4j unrolls the list, because \texttt{b.id IN
\$dsts} planned a \texttt{MultiNodeIndexSeek} that handed the parallel runtime one row (10.3\,s against 0.76\,s at
SF10):
\begin{lstlisting}
CYPHER runtime=parallel MATCH (a:Person {id: $src}) UNWIND $dsts AS d MATCH (b:Person {id: d})
  MATCH p = shortestPath((a)-[:KNOWS*..30]->(b)) RETURN count(p), sum(length(p))
\end{lstlisting}
and Kuzu joins the reached vertices against the destination table:
\begin{lstlisting}
MATCH (a:Person {id: {source}})-[e:Knows* SHORTEST 1..30]->(b:Person)
  WITH b.id AS d, length(e) AS l MATCH (p:benchmark_pairs_same_src_1024 {dst: d}) RETURN count(*), sum(l)
\end{lstlisting}
The whole stream runs in one process per system, without a warm-up. The concurrent driver runs one writer thread and
$R$ reader threads on connections to one \duckdb instance with 32 threads. The writer applies the same batch
statements and pauses 200\,ms after each commit; each reader iteration is one transaction that runs the query on the
index route and on the build route and compares the two results.

\paragraph{Timeouts and measurement.} 600\,s per query. \duckpgq's limit is enforced by killing its process and covers
the whole process, including the untimed warm-up run; Kuzu is interrupted through its connection; Neo4j and Memgraph
transactions are terminated on the server; FalkorDB gets a per-query server timeout of 600\,s. An exceeded budget is
reported, not dropped, and the larger points of the same shape are then recorded as timed out without running them.
Sources $\times$ all is skipped, not timed out, for the native systems when the number of (source, person) pairs
exceeds $5 \times 10^8$. Each point is the mean of three timed runs; the native systems' runs follow each other
without a separate warm-up. \duckpgq's time is the statement time reported by the \duckdb client; the native systems'
times are client-side wall clock, including the driver round trips.
